\section{Mapa documental}

\subsection{Descripcion}

Este documento explica como se organiza el repositorio de documentacion y que pregunta responde cada seccion. Su objetivo es reducir duplicidad, evitar mezclas de responsabilidad y facilitar el mantenimiento.

\subsection{Alcance}

Aplica a todo el repositorio de documentacion tecnica de Quiniela.

\subsection{Definiciones}

\begin{itemize}
  \item \textbf{Seccion}: carpeta principal del repositorio, asociada a un area concreta del sistema.
  \item \textbf{Documento fuente}: archivo `.tex` editable dentro de una subcarpeta `tex/`.
  \item \textbf{Entregable}: PDF publicado en la raiz de una seccion.
  \item \textbf{Documento maestro}: PDF que consolida los documentos aprobados del repositorio.
\end{itemize}

\subsection{Reglas}

\begin{longtable}{p{0.25\textwidth}p{0.67\textwidth}}
\toprule
\textbf{Seccion} & \textbf{Pregunta principal que responde} \\
\midrule
\texttt{00\_Overview} & Que es el sistema, como se organiza la documentacion y bajo que principios debe leerse. \\
\texttt{01\_Producto\_y\_Reglas} & Cuales son las reglas oficiales del producto y del juego. \\
\texttt{02\_Arquitectura} & Como se divide el sistema y cuales son sus limites principales. \\
\texttt{03\_Backend} & Como se organiza el backend a nivel conceptual. \\
\texttt{04\_Frontend} & Como se organiza el frontend a nivel conceptual. \\
\texttt{05\_Datos} & Cual es el modelo de dominio y cuales son sus invariantes. \\
\texttt{06\_Integraciones} & Como interactua el sistema con servicios externos. \\
\texttt{07\_Seguridad\_y\_Operacion} & Que lineamientos operativos y de seguridad deben respetarse. \\
\texttt{08\_Testing} & Como se valida la calidad funcional y tecnica. \\
\texttt{09\_UI\_UX} & Que principios de interfaz y experiencia guian el producto. \\
\texttt{10\_Planeacion} & Que decisiones, riesgos y fases de trabajo siguen abiertas. \\
\bottomrule
\end{longtable}

Una regla normativa debe definirse una sola vez en el repositorio. Los demas documentos deben referenciarla en lugar de duplicarla.

\subsection{Edge Cases}

\begin{itemize}
  \item Si un tema encaja en dos secciones, se documenta en la seccion que lo define y se referencia desde la otra.
  \item Si una regla afecta varias capas del sistema, sigue perteneciendo primero al dominio, no a la capa tecnica que la ejecuta.
\end{itemize}

\subsection{Invariantes}

\begin{itemize}
  \item Cada documento debe tener una sola responsabilidad primaria.
  \item Ninguna regla del negocio debe depender del lector interpretando el contexto.
  \item El documento maestro no reemplaza los documentos modulares; los consolida.
\end{itemize}

\subsection{Preguntas abiertas}

\begin{itemize}
  \item Ninguna por ahora.
\end{itemize}
